%% ==========================================================================
%%  《线性代数：抽象结构 —— 算子及其分解》
%%
%%  编译： xelatex main.tex （建议用 latexmk -xelatex，或在上层目录 make）
%%  正文： 由 tools/md2tex.py 从 Markdown 讲义自动生成，见 body.tex
%% ==========================================================================
\documentclass[UTF8, fontset=none, 11pt, twoside, openright]{ctexbook}

%% ==========================================================================
%%  preamble.tex —— 《线性代数：抽象结构 · 算子及其分解》版式定义
%%  编译引擎：XeLaTeX（需要 ctex + tcolorbox + Libertinus + Noto CJK）
%% ==========================================================================

% -------------------------------------------------------------------------
% 1. 页面尺寸（B5 开本，接近国内数学教材手感）
% -------------------------------------------------------------------------
\usepackage[
  paperwidth  = 176mm,
  paperheight = 250mm,
  inner       = 21mm,
  outer       = 17mm,
  top         = 22mm,
  bottom      = 24mm,
  headheight  = 14pt,
  headsep     = 6mm,
  footskip    = 11mm,
]{geometry}

% -------------------------------------------------------------------------
% 2. 字体：西文 Libertinus，数学 Libertinus Math，中文 Noto Serif/Sans CJK SC
% -------------------------------------------------------------------------
\usepackage{fontspec}
\usepackage{unicode-math}

\setmainfont{LibertinusSerif}[
  Extension      = .otf,
  UprightFont    = *-Regular,
  ItalicFont     = *-Italic,
  BoldFont       = *-Bold,
  BoldItalicFont = *-BoldItalic,
  Numbers        = OldStyle,
  Ligatures      = TeX,
]
\setsansfont{LibertinusSans}[
  Extension      = .otf,
  UprightFont    = *-Regular,
  ItalicFont     = *-Italic,
  BoldFont       = *-Bold,
  Ligatures      = TeX,
]
\setmonofont{LibertinusMono-Regular}[Extension = .otf, Scale = MatchLowercase]
\setmathfont{LibertinusMath-Regular.otf}
% 公式里的数字用等高数字，正文用老式数字
\newfontfamily\liningfont{LibertinusSerif}[Extension=.otf, UprightFont=*-Regular,
  BoldFont=*-Bold, Numbers=Lining]

\xeCJKsetup{CJKmath = true}
\setCJKmainfont{Noto Serif CJK SC}[
  BoldFont       = Noto Serif CJK SC Bold,
  ItalicFont     = Noto Serif CJK SC,
  ItalicFeatures = {FakeSlant = 0.18},
]
\setCJKsansfont{Noto Sans CJK SC}[BoldFont = Noto Sans CJK SC Bold]
\setCJKmonofont{Noto Sans Mono CJK SC}
\newCJKfontfamily\cjkhei{Noto Sans CJK SC}[BoldFont = Noto Sans CJK SC Bold]
\newCJKfontfamily\cjklight{Noto Serif CJK SC}[BoldFont = Noto Serif CJK SC Bold]

% -------------------------------------------------------------------------
% 3. 数学、表格、图形等基础宏包
% -------------------------------------------------------------------------
\usepackage{amsmath, mathtools}   % 注意：unicode-math 已提供全部 AMS 符号，不能再载 amssymb
\usepackage{booktabs, array}
\usepackage{enumitem}
\usepackage{graphicx}
\usepackage{pifont}
\usepackage{microtype}
\usepackage{tikz}
\usetikzlibrary{calc}
\usepackage[most]{tcolorbox}
\tcbuselibrary{theorems, breakable, skins}
\usepackage{titlesec}
\usepackage{titletoc}
\usepackage{fancyhdr}
\usepackage{emptypage}

% -------------------------------------------------------------------------
% 4. 配色方案（沉静的靛蓝为主色，四种辅色区分不同类型的知识块）
% -------------------------------------------------------------------------
\usepackage{xcolor}
\definecolor{inkblue}{HTML}{16324F}      % 主色：章节标题
\definecolor{accent}{HTML}{2C5F8A}       % 辅色：节号、页眉细线
\definecolor{defnblue}{HTML}{1C6E8C}     % 定义
\definecolor{thmred}{HTML}{9B2C2C}       % 定理
\definecolor{propgreen}{HTML}{1F6F54}    % 命题
\definecolor{proofgray}{HTML}{4A5568}    % 证明思路
\definecolor{insightgold}{HTML}{8A6100}  % 要点
\definecolor{ideapurple}{HTML}{5B4B8A}   % 理解
\definecolor{pagegray}{HTML}{6B7280}

% -------------------------------------------------------------------------
% 5. 段落、行距、标点
% -------------------------------------------------------------------------
\linespread{1.30}
\setlength{\parskip}{0.28em plus 0.1em minus 0.05em}
\setlength{\parindent}{2\ccwd}
\setlength{\emergencystretch}{1.5em}
\raggedbottom
\predisplaypenalty=0        % 允许在行间公式前分页，避免大片留白
\postdisplaypenalty=0
\setlength{\abovedisplayskip}{7pt plus 2pt minus 2pt}
\setlength{\belowdisplayskip}{7pt plus 2pt minus 2pt}
\setlength{\abovedisplayshortskip}{3pt plus 2pt}
\setlength{\belowdisplayshortskip}{5pt plus 2pt}

% -------------------------------------------------------------------------
% 6. 超链接与 PDF 书签
% -------------------------------------------------------------------------
\usepackage[
  colorlinks = true,
  linkcolor  = accent,
  citecolor  = accent,
  urlcolor   = accent,
  linktoc    = all,
  bookmarksnumbered  = true,
  bookmarksopen      = true,
  bookmarksopenlevel = 1,
]{hyperref}
\usepackage{bookmark}
% PDF 元信息里含逗号与中文，必须放在 \hypersetup 里而不是宏包选项里
\hypersetup{
  pdftitle    = {线性代数：抽象结构 —— 算子及其分解},
  pdfsubject  = {线性代数讲义},
  pdfkeywords = {线性代数; 算子; 不变子空间; Jordan 标准型; 谱定理; 奇异值分解},
}

% -------------------------------------------------------------------------
% 7. 章 / 节 / 小节标题样式
% -------------------------------------------------------------------------
% 章：深色横条给出章序 + 大标题 + 双色横线
% \CTEXthechapter 与目录中的“第一章”写法保持一致
\providecommand{\CTEXthechapter}{第\,\thechapter\,章}
\titleformat{\chapter}[display]
  {\normalfont\sffamily\raggedright}
  {\colorbox{inkblue}{\parbox{\dimexpr\linewidth-2\fboxsep}{%
     \color{white}\sffamily\bfseries\normalsize
     \strut \CTEXthechapter}}}
  {6pt}
  {\huge\bfseries\color{inkblue}}
  [{\vspace{4pt}{\color{accent}\titlerule[1.1pt]}%
     \vspace{1.4pt}{\color{accent!45}\titlerule[0.5pt]}}]
\titlespacing*{\chapter}{0pt}{4pt}{22pt}

\titleformat{name=\chapter,numberless}[display]
  {\normalfont\sffamily\raggedright}
  {}
  {0pt}
  {\huge\bfseries\color{inkblue}}
  [{\vspace{4pt}{\color{accent}\titlerule[1.1pt]}%
     \vspace{1.4pt}{\color{accent!45}\titlerule[0.5pt]}}]

% 节：色条 + 编号
\titleformat{\section}[hang]
  {\normalfont\sffamily\large\bfseries\color{inkblue}}
  {\raisebox{-0.08em}{\color{accent}\rule{2.6pt}{1.05em}}\hspace{0.55em}%
   \color{accent}\thesection\hspace{0.1em}}
  {0.6em}{}
\titlespacing*{\section}{0pt}{1.7ex plus .3ex}{0.9ex}

\titleformat{\subsection}[hang]
  {\normalfont\sffamily\bfseries\color{accent!85!black}}
  {\thesubsection}{0.55em}{}
\titlespacing*{\subsection}{0pt}{1.3ex plus .2ex}{0.6ex}

% 部分（\part）：整页的大标题，编号用汉字
\renewcommand{\thepart}{\chinese{part}}
\titleformat{\part}[display]
  {\normalfont\sffamily\centering}
  {{\color{accent!45}\rule{0.22\linewidth}{0.9pt}}\\[10pt]
   \color{accent}\Large 第\,\thepart\,部分}
  {14pt}
  {\vspace{2pt}\color{inkblue}\Huge\bfseries}
  [{\vspace{12pt}{\color{accent}\rule{0.42\linewidth}{1.1pt}}}]
\titlespacing*{\part}{0pt}{80pt}{28pt}

% -------------------------------------------------------------------------
% 8. 页眉页脚
% -------------------------------------------------------------------------
\pagestyle{fancy}
\fancyhf{}
\fancyhead[LE]{\small\sffamily\color{pagegray}\leftmark}
\fancyhead[RO]{\small\sffamily\color{pagegray}\rightmark}
\fancyfoot[LE,RO]{\small\sffamily\color{accent}\liningfont\thepage}
\renewcommand{\headrulewidth}{0pt}
\renewcommand{\headrule}{{\color{accent!45}\hrule height 0.4pt}\vspace{1pt}}
\fancypagestyle{plain}{%
  \fancyhf{}%
  \renewcommand{\headrule}{}%
  \fancyfoot[LE,RO]{\small\sffamily\color{accent}\liningfont\thepage}}
\renewcommand{\chaptermark}[1]{\markboth{\CTEXthechapter\quad #1}{}}
\renewcommand{\sectionmark}[1]{\markright{\thesection\quad #1}}

% -------------------------------------------------------------------------
% 9. 目录样式
% -------------------------------------------------------------------------
\renewcommand{\contentsname}{目\hspace{1em}录}
\titlecontents{part}[0pt]
  {\addvspace{18pt}\sffamily\bfseries\color{inkblue}\large}
  {\contentspush{第\,\thecontentslabel\,部分\quad}}{}{}
\titlecontents{chapter}[0pt]
  {\addvspace{7pt}\sffamily\bfseries\color{inkblue}}
  {\contentspush{\thecontentslabel\enspace}}{}
  {\hfill\color{accent}\liningfont\contentspage}
\titlecontents{section}[2.9em]
  {\addvspace{1pt}\small}
  {\contentslabel{2.6em}}{\hspace*{-2.6em}}
  {\titlerule*[0.6pc]{\color{accent!35}.}\color{accent}\liningfont\contentspage}
\titlecontents{subsection}[5.8em]
  {\addvspace{0pt}\small\color{black!80}}
  {\contentslabel{3.1em}}{\hspace*{-3.1em}}
  {\titlerule*[0.6pc]{\color{accent!25}.}\color{accent!80}\liningfont\contentspage}

% -------------------------------------------------------------------------
% 10. 知识块：定义 / 定理 / 命题 / 证明思路 / 要点 / 理解 / 部分导读
% -------------------------------------------------------------------------
% 盒子内的段落：不缩进，靠空白分段（parbox=false 才能让段设置生效）
\newcommand{\jboxpar}{\setlength{\parindent}{0pt}\setlength{\parskip}{0.5em plus 0.1em}}

% 统一外观：隐去边框，只留左侧色条 + 极淡底色；三类共用一个“结果”计数器
\tcbset{
  jstatement/.style = {
    enhanced jigsaw, breakable, sharp corners, boxrule = 0pt, frame hidden,
    parbox = false,
    left = 10pt, right = 10pt, top = 7pt, bottom = 7pt,
    before skip = 11pt plus 2pt, after skip = 11pt plus 2pt,
    fonttitle = \sffamily\bfseries,
    theorem style = standard,
    separator sign none,
    terminator sign none,
    description delimiters = {\kern-0.42em（}{）},
    before upper = {\jboxpar},
  },
}

\newtcbtheorem[number within = chapter, list inside = jdeflist]{jdefn}{定义}{
  jstatement,
  colback = defnblue!5!white,
  coltitle = defnblue!88!black,
  borderline west = {2.6pt}{0pt}{defnblue},
}{def}

\newtcbtheorem[use counter from = jdefn, list inside = jthmlist]{jthm}{定理}{
  jstatement,
  colback = thmred!5!white,
  coltitle = thmred!88!black,
  borderline west = {2.6pt}{0pt}{thmred},
}{thm}

\newtcbtheorem[use counter from = jdefn, list inside = jproplist]{jprop}{命题}{
  jstatement,
  colback = propgreen!5!white,
  coltitle = propgreen!88!black,
  borderline west = {2.6pt}{0pt}{propgreen},
}{prop}

% 证明思路：无底色，细灰竖线，末尾一个小方块（用 \rule 画，不依赖字体字形）
\newcommand{\jqed}{\textcolor{proofgray!80}{\rule[0.03em]{0.46em}{0.46em}}}
\newtcolorbox{jproof}[1]{
  enhanced jigsaw, breakable, sharp corners, boxrule = 0pt, frame hidden,
  parbox = false,
  colback = white,
  borderline west = {1pt}{0pt}{proofgray!45},
  left = 9pt, right = 6pt, top = 4pt, bottom = 4pt,
  before skip = 9pt plus 2pt, after skip = 11pt plus 2pt,
  fontupper = \small,
  before upper = {\jboxpar\textcolor{proofgray}{\sffamily\bfseries #1\,}\hspace{0.45em}},
  after upper  = {\unskip\nobreak\hfill\jqed},
}

% 要点（原 Markdown 的引用块）
\newtcolorbox{jinsight}{
  enhanced jigsaw, breakable, sharp corners, boxrule = 0pt, frame hidden,
  parbox = false,
  colback = insightgold!8!white,
  borderline west = {2.6pt}{0pt}{insightgold!85},
  left = 10pt, right = 10pt, top = 7pt, bottom = 7pt,
  before skip = 11pt plus 2pt, after skip = 11pt plus 2pt,
  before upper = {\jboxpar\textcolor{insightgold!85!black}{\ding{72}\ \sffamily\bfseries 要点\,}\hspace{0.4em}},
}

% 理解 / 构造目标一类的旁注
\newtcolorbox{jidea}[1]{
  enhanced jigsaw, breakable, sharp corners, boxrule = 0pt, frame hidden,
  parbox = false,
  colback = ideapurple!5!white,
  borderline west = {2.6pt}{0pt}{ideapurple!85},
  left = 10pt, right = 10pt, top = 6pt, bottom = 6pt,
  before skip = 10pt plus 2pt, after skip = 10pt plus 2pt,
  fontupper = \small,
  before upper = {\jboxpar\textcolor{ideapurple!85!black}{\sffamily\bfseries #1\,}\hspace{0.45em}},
}

% 每一“部分”的导读：整页宽的浅色大盒子
\newtcolorbox{partintro}[1]{
  enhanced, breakable, sharp corners,
  parbox = false,
  colback = accent!4!white, colframe = accent!35,
  boxrule = 0.6pt,
  borderline west = {3pt}{0pt}{accent},
  left = 14pt, right = 14pt, top = 14pt, bottom = 12pt,
  before skip = 6pt, after skip = 16pt,
  fonttitle = \sffamily\bfseries,
  coltitle = white,
  colbacktitle = accent,
  title = {\ #1},
  before upper = {\jboxpar},
  attach boxed title to top left = {xshift = 6pt, yshift = -\tcboxedtitleheight/2},
  boxed title style = {sharp corners, boxrule = 0pt, colback = accent},
}

% -------------------------------------------------------------------------
% 11. 列表样式
% -------------------------------------------------------------------------
\newcommand{\jsquare}{\textcolor{accent}{\raisebox{0.22em}{\rule{0.42em}{0.42em}}}}
\newlist{jitem}{itemize}{3}
\setlist[jitem]{label = \jsquare,
  leftmargin = 1.9em, itemsep = 1pt, topsep = 4pt, parsep = 2pt}
\newlist{jenum}{enumerate}{3}
\setlist[jenum]{label = {\textcolor{accent}{\sffamily\bfseries\arabic*.}},
  leftmargin = 2.1em, itemsep = 1pt, topsep = 4pt, parsep = 2pt}

% -------------------------------------------------------------------------
% 12. 常用数学记号
% -------------------------------------------------------------------------
\newcommand{\FF}{\mathbb{F}}
\newcommand{\RR}{\mathbb{R}}
\newcommand{\CC}{\mathbb{C}}
\DeclareMathOperator{\spann}{span}
\DeclareMathOperator{\diag}{diag}

% 扉页上用的装饰线
\newcommand{\titlerulebig}[1]{{\color{#1}\rule{\linewidth}{1.2pt}}}


% ---- 书名、作者等元信息（改这里即可）-------------------------------------
\newcommand{\bookseries}{线性代数讲义\; \textbullet\; 抽象结构}
\newcommand{\booktitleA}{算子}
\newcommand{\booktitleB}{及其分解}
\newcommand{\authorone}{夜航未西飞}
\newcommand{\authoronerole}{著}
\newcommand{\authortwo}{阿酒}
\newcommand{\authortworole}{校对\;\textbullet\;\LaTeX{} 排版}
\newcommand{\bookdate}{\the\year 年\ \the\month 月}

\begin{document}

% =========================================================================
%  扉页
% =========================================================================
\begin{titlepage}
\thispagestyle{empty}
\begin{tikzpicture}[remember picture, overlay]
  % 左侧竖向色带（宽带 + 一条浅色细带）
  \fill[mainblue] (current page.north west) rectangle
    ($(current page.south west) + (16mm, 0)$);
  \fill[mainblue!42!white] ($(current page.north west) + (16mm, 0)$) rectangle
    ($(current page.south west) + (19mm, 0)$);
  % 右上角小方块装饰
  \fill[mainblue!30!white] ($(current page.north east) + (-30mm, -18mm)$) rectangle
    ($(current page.north east) + (-24mm, -24mm)$);
  \fill[thmred] ($(current page.north east) + (-22mm, -18mm)$) rectangle
    ($(current page.north east) + (-16mm, -24mm)$);
\end{tikzpicture}%
\vspace*{38mm}

\begin{flushleft}
  {\sffamily\large\color{accent}\bookseries}\\[2mm]
  {\color{accent!45}\rule{0.40\linewidth}{0.9pt}}\\[10mm]
  {\sffamily\bfseries\color{inkblue}\fontsize{46}{52}\selectfont \booktitleA}\\[4mm]
  {\sffamily\bfseries\color{inkblue}\fontsize{32}{38}\selectfont \booktitleB}\\[12mm]
  {\color{inkblue}\rule{0.62\linewidth}{1.4pt}}\\[7mm]
  {\large\color{black!72}
   不变子空间 \quad\textbullet\quad 上三角化与对角化\\[2.4mm]
   Jordan 标准型 \quad\textbullet\quad 实向量空间的复化\\[2.4mm]
   伴随与谱定理 \quad\textbullet\quad 极分解与奇异值分解}
\end{flushleft}

\vfill
\begin{flushleft}
  {\sffamily\bfseries\Large\color{mainblue}\authorone}%
  \hspace{0.9em}{\small\color{pagegray}\authoronerole}\\[3.2mm]
  {\sffamily\large\color{mainblue!85!black}\authortwo}%
  \hspace{0.9em}{\small\color{pagegray}\authortworole}\\[6mm]
  {\sffamily\color{pagegray}\bookdate}
\end{flushleft}
\end{titlepage}

% =========================================================================
%  版权/说明页
% =========================================================================
\thispagestyle{empty}
\null\vfill
\begin{flushleft}
\small\color{black!75}
\textsf{\textbf{\color{inkblue}关于这本讲义}}\\[4pt]
本书的主题是有限维向量空间上的\emph{算子}以及它的各种分解：
从不变子空间出发，经由特征结构、上三角化、对角化、Jordan 标准型与实向量空间的复化，
再进入内积空间，讨论伴随、正规算子、自伴算子、谱定理，最后到极分解与奇异值分解。\\[10pt]

\textsf{\textbf{\color{inkblue}分工}}\\[4pt]
\authorone{}：讲义正文的构思与撰写。\\
\authortwo{}：后期校对、\LaTeX{} 排版与构建流程。\\[10pt]

\textsf{\textbf{\color{inkblue}排版说明}}\\[4pt]
正文由 Markdown 源稿经 \texttt{tools/md2tex.py} 自动转换为 \LaTeX{}，
用 XeLaTeX 编译；西文字体为 Libertinus，公式字体为 Libertinus Math，
中文字体为思源宋体 / 思源黑体（Noto Serif~/~Sans CJK SC）。\\[10pt]

\textsf{\textbf{\color{inkblue}版本}}\\[4pt]
\bookdate 编译。源码与修订历史由 Git 管理。
\end{flushleft}
\vspace*{18mm}

% =========================================================================
%  前置部分：体例说明 + 目录
% =========================================================================
\frontmatter
\pagestyle{fancy}

\chapter*{阅读指引：本书的体例}
\addcontentsline{toc}{chapter}{阅读指引：本书的体例}
\markboth{阅读指引}{阅读指引}

全书用颜色区分四类内容，正文中的每一个色块都可以单独阅读：

\begin{jdefn*}{定义框}{}
藏青蓝色块给出概念的定义，以及为了让定义站得住脚所需要的最少解释。
\end{jdefn*}

\begin{jthm*}{定理框}{}
暗红色块给出定理。定理只陈述结论；论证放在紧随其后的灰色证明块里。
\end{jthm*}

\begin{jprop*}{命题框}{}
墨绿色块给出性质、刻画与推论——它们通常是定理的零件，或定义的直接后果。
\end{jprop*}

\begin{jproof}{证明思路}
灰色竖线标出的段落是证明或证明思路。本书里许多证明只给思路：
先说清楚“为什么必须这样做”，再补细节。段末的小方块表示论证结束。
\end{jproof}

\begin{jidea}{理解}
靖紫色块是对刚才那段形式化内容的翻译，回答“这件事到底在说什么”。
\end{jidea}

\begin{jinsight}
琥珀金色块是一章或一节的要点：如果只记一句话，就记这一句。
\end{jinsight}

未加色块的部分是叙述性的正文，负责把这些结论串成一条线。
每一“部分”的开头都有一页藏青蓝导读，说明该部分研究什么、用到哪些结构。

\cleardoublepage
\tableofcontents

% =========================================================================
%  正文
% =========================================================================
\mainmatter
%% body.tex —— 全书骨架（由 tools/md2tex.py 自动生成）

%% 绪论：算子与它的分解
%% 本文件由 tools/md2tex.py 自动生成，请勿直接编辑。

\chapter*{绪论\quad 算子与它的分解}
\addcontentsline{toc}{chapter}{绪论\quad 算子与它的分解}
\markboth{绪论}{绪论}


\section*{这一部分研究什么}
\addcontentsline{toc}{section}{这一部分研究什么}

算子是从一个向量空间映到它自身的线性映射：

\[
T:V\to V.
\]

因为输入与输出都在同一空间中，算子可以反复作用：

\[
v,\quad T(v),\quad T^2(v),\quad T^3(v),\ldots
\]

因此，算子不仅描述一次线性变化，也描述一个线性系统内部不断演化的结构。

本章研究的核心问题是：

\begin{jinsight}
能否将一个复杂算子拆成若干简单、彼此独立或具有明确耦合方式的部分？
\end{jinsight}

若可以找到一组适合的基，使算子的矩阵具有对角、块对角或 Jordan 块等简单形状，我们就能直接读出算子的长期作用、可逆性、特征值、幂的行为与不变结构。这对于物理系统的分析，微分方程的解，机器学习中权重矩阵的分析，研究线性方程组的解都是非常有用的。

\section*{为什么可以分解}
\addcontentsline{toc}{section}{为什么可以分解}

算子分解的共同基础是不变子空间。

若 $U$ 是 $V$ 的子空间，且：

\[
T(U)\subseteq U,
\]

则称 $U$ 是 $T$ 的不变子空间。

这意味着：一个向量只要从 $U$ 出发，无论施加 $T$ 多少次，都不会离开 $U$。

若存在：

\[
V=U_1\oplus\cdots\oplus U_m,
\]

并且每个 $U_j$ 都对 $T$ 不变，那么任意向量都能唯一写成：

\[
v=u_1+\cdots+u_m,
\qquad u_j\in U_j.
\]

由线性性：

\[
T(v)=T(u_1)+\cdots+T(u_m).
\]

每一部分仍留在自己的不变子空间中。因此，我们可以分别研究 $T$ 在每个 $U_j$ 上的作用，再将结果组合起来。

算子分解的本质，就是寻找足够多的不变子空间，并理解它们之间是否能够形成直和。

\section*{分解方式总览}
\addcontentsline{toc}{section}{分解方式总览}


\subsection*{一般实向量空间与复向量空间}
\addcontentsline{toc}{subsection}{一般实向量空间与复向量空间}

这一部分只使用线性结构，不要求内积。

\begin{jitem}
\item %
上三角化：将算子表示为沿一条不变子空间链逐层作用的形式；
\item %
对角化：将空间分解为一维不变子空间的直和；
\item %
Jordan 标准型：当特征向量不够时，用广义特征向量和 Jordan 链描述剩余耦合；
\item %
复化：将实算子扩展到复空间中，以暴露复特征值与复 Jordan 结构，再翻译回实空间中的二维不变块。
\end{jitem}

其中，对角化是 Jordan 标准型的特殊情形：每一个 Jordan 块都只有 $1\times1$ 时，Jordan 矩阵就是对角矩阵。

\subsection*{内积空间中的算子分解}
\addcontentsline{toc}{subsection}{内积空间中的算子分解}

内积为向量空间增加长度、正交与伴随结构。

对于自伴算子、正规算子、酉算子等与内积相容的特殊算子，谱定理能够给出规范正交的特征向量基。此时算子不仅可对角化，而且可以在保持长度和正交关系的坐标中对角化。

\subsection*{一般线性映射的分解}
\addcontentsline{toc}{subsection}{一般线性映射的分解}

极分解与奇异值分解更进一步：它们不要求映射的定义域与陪域相同。

若：

\[
T:V\to W,
\]

其中 $V,W$ 是有限维内积空间，则奇异值分解可以在输入空间和输出空间中分别选取规范正交基，使 $T$ 只沿对应方向进行缩放。

极分解则将 $T$ 写成：

\[
T=U|T|,
\qquad
|T|=(T^*T)^{1/2}.
\]

其中 $|T|$ 是定义在输入空间上的正半定自伴算子；$U$ 一般是部分等距算子。只有在适当的单射、满射或可逆条件下，$U$ 才会成为真正的等距同构。

\section*{本章的结构}
\addcontentsline{toc}{section}{本章的结构}

本章先讨论一般实、复向量空间上的算子：不变子空间、特征结构、上三角化、对角化、Jordan 标准型与复化。

之后再讨论内积空间：伴随、正规算子、自伴算子、谱定理、极分解与奇异值分解。

%% 第 1 部分：一般实向量空间与复向量空间上的算子分解
%% 本文件由 tools/md2tex.py 自动生成，请勿直接编辑。

\part{一般实向量空间与复向量空间上的算子分解}

\thispagestyle{plain}
\begin{partintro}{这一部分研究什么}
这一部分只使用向量空间与线性映射的结构，不使用与内积相关的内容。

设：

\[
T:V\to V
\]

是一个有限维向量空间上的算子。

我们的目标是寻找合适的基，或将 $V$ 分解为若干个对 $T$ 不变的子空间，使 $T$ 的作用变得简单、可读。

在复向量空间上，最终可以将任意算子写成 Jordan 标准型；对角化是 Jordan 标准型最简单的特殊情形。

在实向量空间上，某些算子没有实特征值。此时需要先复化，观察其在复空间中的 Jordan 结构，再将共轭成对的复结构翻译回实空间中的二维不变块。
\end{partintro}

%% 第 1 章：算子与不变子空间
%% 本文件由 tools/md2tex.py 自动生成，请勿直接编辑。

\chapter{算子与不变子空间}


\begin{jdefn}{算子}{r:2:1}
从向量空间 $V$ 映到自身的线性映射：

\[
T:V\to V
\]

称为 $V$ 上的算子。

因为定义域与陪域相同，算子可以反复作用：

\[
v,\quad T(v),\quad T^2(v),\quad T^3(v),\ldots
\]

其中：

\[
T^0=I,
\qquad
T^{k+1}=T\circ T^k.
\]
\end{jdefn}

\begin{jdefn}{不变子空间}{r:2:2}
若 $U\le V$，且：

\[
T(U)\subseteq U,
\]

则称 $U$ 是 $T$ 的不变子空间。

这表示：从 $U$ 中任取一个向量，反复施加 $T$ 后，结果始终不会离开 $U$。
\end{jdefn}

\section{不变子空间为何是分解的基础}

若：

\[
V=U_1\oplus\cdots\oplus U_m,
\]

并且每个 $U_j$ 都对 $T$ 不变，则任意向量可唯一写成：

\[
v=u_1+\cdots+u_m,
\qquad
u_j\in U_j.
\]

由线性性：

\[
T(v)=T(u_1)+\cdots+T(u_m).
\]

而每个 $T(u_j)$ 仍位于对应的 $U_j$ 中。

因此，算子在不同不变子空间上的作用彼此不会混合。我们可以先分别研究 $T|_{U_j}$，再把结论合并。

\begin{jinsight}
算子分解的核心，就是寻找足够多的不变子空间，并判断它们能否构成直和。
\end{jinsight}

%% 第 2 章：限制算子与商算子
%% 本文件由 tools/md2tex.py 自动生成，请勿直接编辑。

\chapter{限制算子与商算子}


\begin{jdefn}{限制算子}{r:3:1}
若 $U$ 是 $T$ 的不变子空间，则可以只研究 $T$ 在 $U$ 内部的作用：

\[
T|_U:U\to U.
\]

它称为 $T$ 在 $U$ 上的限制算子。

限制算子用于研究已经被分离出的部分。
\end{jdefn}

\begin{jdefn}{商算子}{r:3:2}
若 $U$ 是 $T$ 的不变子空间，则可在商空间 $V/U$ 上定义：

\[
\widetilde T:V/U\to V/U,
\]

\[
\widetilde T(v+U)=T(v)+U.
\]

它称为 $T$ 关于 $U$ 的商算子。
\end{jdefn}

\section{商算子为何良定义}

若：

\[
v+U=w+U,
\]

则：

\[
v-w\in U.
\]

因为 $U$ 对 $T$ 不变：

\[
T(v-w)\in U.
\]

于是：

\[
T(v)+U=T(w)+U.
\]

所以 $\widetilde T(v+U)$ 不依赖陪集代表元的选择，是良定义的。

\section{限制与商的意义}

\begin{jitem}
\item %
限制算子研究已经找到的不变部分；
\item %
商算子研究忽略该部分后，剩余自由度上的作用。
\end{jitem}

它们是上三角化、主分解与 Jordan 标准型归纳构造的重要工具。

%% 第 3 章：特征值、特征向量与特征空间
%% 本文件由 tools/md2tex.py 自动生成，请勿直接编辑。

\chapter{特征值、特征向量与特征空间}


\begin{jdefn}{特征值与特征向量}{r:4:1}
若存在非零向量 $v\in V$，使：

\[
T(v)=\lambda v,
\]

则称 $\lambda\in\mathbb F$ 是 $T$ 的特征值，$v$ 是对应的特征向量。

特征向量张成的一维空间：

\[
\operatorname{span}(v)
\]

是最简单的不变子空间。
\end{jdefn}

\section{特征值的等价描述}

\[
\lambda\text{ 是 }T\text{ 的特征值}
\iff
\ker(T-\lambda I)\ne\{0\}.
\]

在有限维情形，还等价于：

\[
T-\lambda I\text{ 不可逆}.
\]

注意：

\[
T-\lambda I=0
\]

表示 $T=\lambda I$，即空间中每个向量都被同一个比例缩放。这比“$\lambda$ 是一个特征值”强得多。

\begin{jdefn}{特征空间}{r:4:2}
定义：

\[
E_\lambda
=
\ker(T-\lambda I).
\]

它称为 $\lambda$ 对应的特征空间。

其中的非零向量正是所有属于 $\lambda$ 的特征向量。
\end{jdefn}

\begin{jprop}{不同特征值对应的特征向量线性无关}{r:4:3}
若：

\[
\lambda_1,\ldots,\lambda_m
\]

两两不同，且：

\[
v_j\in E_{\lambda_j}\setminus\{0\},
\]

则：

\[
v_1,\ldots,v_m
\]

线性无关。
\end{jprop}

\begin{jproof}{证明思路}
假设它们存在最短的非平凡线性关系。

对该关系施加 $T-\lambda_m I$。最后一个向量会被消掉，而其余特征向量只会被乘上非零标量，因此得到一条更短的非平凡关系，产生矛盾。
\end{jproof}

\section{同一特征值也可能有多个线性无关特征向量}

若：

\[
\dim E_\lambda>1,
\]

则同一个特征值 $\lambda$ 对应多个线性无关特征向量。

%% 第 4 章：算子多项式
%% 本文件由 tools/md2tex.py 自动生成，请勿直接编辑。

\chapter{算子多项式}


\begin{jdefn}{算子多项式的定义}{r:5:1}
对于多项式：

\[
p(z)=a_0+a_1z+\cdots+a_mz^m,
\]

定义：

\[
p(T)
=
a_0I+a_1T+\cdots+a_mT^m.
\]

这个定义有意义，因为 $I,T,T^2,\ldots$ 都是从 $V$ 到 $V$ 的线性算子，可以相加与数乘。
\end{jdefn}

\begin{jprop}{算子多项式的乘法}{r:5:2}
若 $p,q$ 是两个多项式，则：

\[
p(T)q(T)=(pq)(T).
\]
\end{jprop}

\begin{jidea}{理解}
所有出现在 $p(T)$ 与 $q(T)$ 中的项，都是同一个算子 $T$ 的幂。因此：

\[
T^rT^s=T^{r+s},
\]

且它们彼此可交换。

这使普通多项式的乘法规则可以直接迁移到算子多项式。
\end{jidea}

\begin{jprop}{多项式算子的核和值域是不变子空间}{r:5:3}
对于任意多项式 $p$：

\[
\ker p(T)
\]

与：

\[
\operatorname{range}p(T)
\]

都是 $T$ 的不变子空间。
\end{jprop}

\begin{jproof}{证明思路}
$T$ 与 $p(T)$ 可交换。

若 $v\in\ker p(T)$，则：

\[
p(T)T(v)=T p(T)(v)=0.
\]

所以 $T(v)$ 仍在 $\ker p(T)$ 中。

若 $v\in\operatorname{range}p(T)$，则存在 $u$ 使：

\[
v=p(T)u.
\]

于是：

\[
T(v)=Tp(T)u=p(T)T(u).
\]

所以 $T(v)$ 仍在 $\operatorname{range}p(T)$ 中。
\end{jproof}

%% 第 5 章：代数基本定理及其在复空间算子上的推论
%% 本文件由 tools/md2tex.py 自动生成，请勿直接编辑。

\chapter{代数基本定理及其在复空间算子上的推论}


\begin{jthm}{代数基本定理}{r:6:1}
任意非常数复系数多项式：

\[
p(z)=a_nz^n+\cdots+a_1z+a_0,
\qquad
a_n\ne0,
\]

至少有一个复根。

等价地，任意 $n$ 次复系数多项式都可以完全分解为一次因子的乘积：

\[
p(z)
=
a_n(z-\lambda_1)\cdots(z-\lambda_n).
\]

其中 $\lambda_1,\ldots,\lambda_n$ 允许重复。重复次数称为相应根的重数。
\end{jthm}

\section{为什么它对算子重要}

设 $V$ 是一个非零有限维复向量空间，且：

\[
T:V\to V.
\]

$T$ 的特征多项式：

\[
p_T(z)=\det(zI-T)
\]

是一个次数为：

\[
\dim V
\]

的复系数多项式。

因此，由代数基本定理，$p_T$ 至少有一个复根 $\lambda$：

\[
p_T(\lambda)=0.
\]

这意味着：

\[
\det(\lambda I-T)=0.
\]

所以：

\[
\lambda I-T
\]

不可逆；等价地：

\[
T-\lambda I
\]

的零空间非平凡。

因此，存在非零向量 $v$ 满足：

\[
(T-\lambda I)v=0.
\]

也就是：

\[
T(v)=\lambda v.
\]

所以得到结论：

\begin{jthm}{复向量空间上的算子一定存在特征值}{r:6:2}
若 $V$ 是非零有限维复向量空间，则任意算子：

\[
T:V\to V
\]

至少存在一个特征值。
\end{jthm}

\section{这个推论的意义}

这保证了复空间上的任意算子至少存在一个一维不变子空间：

\[
\operatorname{span}(v).
\]

因此，我们总能先从一个特征向量出发，再通过限制算子、商算子或归纳过程继续分析其余空间。

这正是复数域上任意算子都可以上三角化的起点。

%% 第 6 章：特征多项式、Cayley-Hamilton 定理与最小多项式
%% 本文件由 tools/md2tex.py 自动生成，请勿直接编辑。

\chapter{特征多项式、Cayley-Hamilton 定理与最小多项式}


\begin{jdefn}{特征多项式}{r:7:1}
设 $V$ 有限维。

定义 $T$ 的特征多项式为：

\[
p_T(z)=\det(zI-T).
\]

其中，$I$ 是恒等算子在某组基下的单位矩阵表示。
\end{jdefn}

\begin{jprop}{特征值与特征多项式}{r:7:2}
\[
\lambda\text{ 是 }T\text{ 的特征值}
\iff
p_T(\lambda)=0.
\]
\end{jprop}

\begin{jidea}{理解}
\[
p_T(\lambda)=0
\]

意味着：

\[
\det(\lambda I-T)=0.
\]

因此 $\lambda I-T$ 不可逆，也就是：

\[
T-\lambda I
\]

不具有可逆性。于是其零空间非平凡，存在特征向量。
\end{jidea}

\begin{jdefn}{代数重数与几何重数}{r:7:3}
若 $\lambda$ 是 $p_T(z)$ 的根，则它作为根的重数称为 $\lambda$ 的代数重数。

\[
\dim E_\lambda
\]

称为 $\lambda$ 的几何重数。

几何重数表示有多少个独立的普通特征向量；代数重数则表示该特征值在特征多项式中总共出现多少次。
\end{jdefn}

\begin{jthm}{Cayley-Hamilton 定理}{r:7:4}
\[
p_T(T)=0.
\]

即：每个有限维算子都满足自己的特征多项式方程。
\end{jthm}

\section{Cayley-Hamilton 定理有什么用}

在复数域上，代数基本定理保证：

\[
p_T(z)
=
\prod_{j=1}^s(z-\lambda_j)^{a_j},
\]

其中：

\[
\lambda_1,\ldots,\lambda_s
\]

是全部不同特征值。

由 Cayley-Hamilton 定理：

\[
\prod_{j=1}^s(T-\lambda_jI)^{a_j}=0.
\]

这说明：反复施加各个 $T-\lambda_jI$ 后，整个空间中的任何向量最终都会被消掉。这个定理的证明从略。

不同特征值对应的因子彼此互素。正是这种互素性，使我们能够把空间拆成不同特征值对应的广义特征空间。

\begin{jinsight}
代数基本定理负责把特征多项式拆成不同特征值的因子；Cayley-Hamilton 定理则保证这些因子写成算子的形式之后，真的能把整个空间分解成广义特征空间。
\end{jinsight}

\section{不同广义特征空间为何能够构成直和}

设特征多项式分解为：

\[
p_T(z)
=
\prod_{j=1}^s(z-\lambda_j)^{a_j},
\]

其中：

\[
\lambda_1,\ldots,\lambda_s
\]

两两不同。

由 Cayley-Hamilton 定理：

\[
\prod_{j=1}^s(T-\lambda_jI)^{a_j}=0.
\]

这里仅说明：全部因子依次作用后会消掉任意向量。要将空间拆成不同特征值对应的部分，还需要利用不同因子之间的互素性。

以两个因子为例，令：

\[
f(z)=(z-\lambda)^a,
\qquad
g(z)=(z-\mu)^b,
\qquad
\lambda\ne\mu.
\]

因为 $f$ 与 $g$ 没有公共根，所以它们互素。由 Bézout 恒等式，存在多项式 $r,s$，使：

\[
r(z)f(z)+s(z)g(z)=1.
\]

将 $z$ 替换为 $T$，得到：

\[
r(T)f(T)+s(T)g(T)=I.
\]

因此，对任意 $v\in V$：

\[
v
=
r(T)f(T)v+s(T)g(T)v.
\]

其中，第一项被 $g(T)$ 消掉，第二项被 $f(T)$ 消掉。因此：

\[
v
\in
\ker f(T)+\ker g(T).
\]

若 $v$ 同时属于：

\[
\ker f(T)
\quad\text{与}\quad
\ker g(T),
\]

则：

\[
v
=
r(T)f(T)v+s(T)g(T)v
=
0.
\]

所以：

\[
\ker f(T)\cap\ker g(T)=\{0\}.
\]

故：

\[
V
=
\ker f(T)\oplus\ker g(T).
\]

对全部两两不同的特征值重复这一论证，即得到：

\[
V
=
G_{\lambda_1}\oplus\cdots\oplus G_{\lambda_s},
\]

其中：

\[
G_{\lambda_j}
=
\ker(T-\lambda_jI)^{a_j}.
\]

\begin{jinsight}
这是由代数性质证明的，而在下面会有另外一种方式来证明，$V$ 可以被直和分解，而它依赖于我们之后证明的有关幂零算子的性质。
\end{jinsight}

\begin{jdefn}{最小多项式}{r:7:5}
使：

\[
m_T(T)=0
\]

成立的所有首一多项式中，次数最小的那个称为 $T$ 的最小多项式，记作：

\[
m_T(z).
\]
\end{jdefn}

\begin{jprop}{最小多项式的性质}{r:7:6}
\[
m_T(z)\mid p_T(z).
\]

并且，任何满足：

\[
q(T)=0
\]

的多项式 $q$，都能被 $m_T$ 整除。
\end{jprop}

\section{最小多项式的意义}

特征多项式告诉我们每个特征值总共出现多少次。

最小多项式告诉我们：对每个特征值，算子偏离可对角化状态的最深程度，即最大需要应用几次幂零算子才能形成一条完整的若尔当链。

在 Jordan 标准型中，若：

\[
m_T(z)
=
\prod_\lambda(z-\lambda)^{r_\lambda},
\]

则 $r_\lambda$ 等于 $\lambda$ 对应的最大 Jordan 块的大小。

%% 第 7 章：复向量空间上的上三角化
%% 本文件由 tools/md2tex.py 自动生成，请勿直接编辑。

\chapter{复向量空间上的上三角化}


\begin{jthm}{上三角化定理}{r:8:1}
若 $V$ 是有限维复向量空间，则任意算子：

\[
T:V\to V
\]

都至少存在一组基，使其矩阵为上三角矩阵。
\end{jthm}

\section{为什么只有复数域下总能上三角化}

证明的起点是代数基本定理。

$T$ 的特征多项式是一个复系数多项式，因此至少有一个复根。也就是说，$T$ 至少有一个特征值与一个一维不变子空间。

之后将该一维不变子空间取商，在商空间上继续寻找特征值。不断重复，便得到一条逐层增长的不变子空间链。这里的证明从略。

\begin{jthm}{舒尔定理}{r:8:2}
在复内积空间下，则总存在至少一组规范正交基使其矩阵为上三角矩阵，这称为舒尔定理
\end{jthm}

\section{上三角矩阵与不变子空间链}

若存在：

\[
\{0\}\subset V_1\subset V_2\subset\cdots\subset V_n=V,
\]

其中：

\[
\dim V_j=j,
\]

且每个 $V_j$ 对 $T$ 不变，那么可以选取基：

\[
(v_1,\ldots,v_n),
\]

使：

\[
V_j=\operatorname{span}(v_1,\ldots,v_j).
\]

于是：

\[
T(v_j)\in\operatorname{span}(v_1,\ldots,v_j).
\]

因此，$T(v_j)$ 的坐标在第 $j$ 行以下都为零。由于矩阵第 $j$ 列记录 $T(v_j)$，矩阵必为上三角形式。

\begin{jinsight}
上三角矩阵的本质，是一条逐层嵌套的不变子空间链。
\end{jinsight}

\begin{jprop}{上三角矩阵的可逆性}{r:8:3}
对于一个算子 $T$ 在某组基下的上三角矩阵，它是可逆的当且仅当全部对角线元素非零

证明思路：若存在对角线 $0$ 元素，则意味着，存在多个向量能够被映射为同一个结果，这意味着映射并非单射，因此零空间非零。
\end{jprop}

\begin{jprop}{上三角矩阵的特征值}{r:8:4}
若 $T$ 的上三角矩阵对角元为：

\[
a_1,\ldots,a_n,
\]

则这些对角元就是 $T$ 的全部特征值，计入代数重数。
\end{jprop}

\begin{jproof}{证明思路}
\[
T-\lambda I
\]

仍然是上三角矩阵，其对角元是：

\[
a_1-\lambda,\ldots,a_n-\lambda.
\]

上三角矩阵可逆，当且仅当所有对角元都非零。

因此，$T-\lambda I$ 不可逆，当且仅当：

\[
\lambda=a_j
\]

对某个 $j$ 成立。
\end{jproof}

%% 第 8 章：对角化
%% 本文件由 tools/md2tex.py 自动生成，请勿直接编辑。

\chapter{对角化}


\begin{jdefn}{对角化的定义}{r:9:1}
算子 $T$ 可对角化，当且仅当存在一组由特征向量构成的基。

在该基下：

\[
[T]
=
\begin{pmatrix}
\lambda_1&0&\cdots&0\\
0&\lambda_2&\cdots&0\\
\vdots&\vdots&\ddots&\vdots\\
0&0&\cdots&\lambda_n
\end{pmatrix}.
\]
\end{jdefn}

\section{对角化的结构意义}

对角化等价于：

\[
V
=
\bigoplus_\lambda E_\lambda,
\]

其中求和遍历 $T$ 的全部不同特征值。

每个特征空间本身可能高于一维；但在特征空间内部，$T$ 只做同一个比例的缩放：

\[
T(v)=\lambda v,
\qquad
v\in E_\lambda.
\]

\section{对角化的意义}

在特征向量基下，算子不再混合不同方向。

对每个坐标方向，$T$ 只进行独立缩放。于是：

\[
T^k(v_j)=\lambda_j^k v_j.
\]

这使算子的高次幂、线性递推与长期行为都变得直接可读。

\section{对角化与 Jordan 标准型}

对角化是 Jordan 标准型的特殊情形。

当每个 Jordan 块都只有：

\[
1\times1
\]

时，Jordan 标准型就是一个对角矩阵。

%% 第 9 章：幂零算子与零空间链
%% 本文件由 tools/md2tex.py 自动生成，请勿直接编辑。

\chapter{幂零算子与零空间链}


\begin{jdefn}{幂零算子}{r:10:1}
若存在正整数 $q$，使：

\[
N^q=0,
\]

则称 $N$ 是幂零算子。

幂零算子反复作用足够多次后，会将整个空间中的所有向量映为零。
\end{jdefn}

\begin{jprop}{零空间链递增}{r:10:2}
对任意算子 $N$：

\[
\ker N
\subseteq
\ker N^2
\subseteq
\ker N^3
\subseteq\cdots.
\]
\end{jprop}

\begin{jproof}{证明思路}
若：

\[
N^k(v)=0,
\]

则再施加一次 $N$ 后仍为零：

\[
N^{k+1}(v)=0.
\]

因此：

\[
v\in\ker N^{k+1}.
\]
\end{jproof}

\begin{jprop}{零空间链稳定}{r:10:3}
若：

\[
\ker N^r=\ker N^{r+1},
\]

则：

\[
\ker N^r=\ker N^{r+1}=\ker N^{r+2}=\cdots.
\]
\end{jprop}

\begin{jproof}{证明思路}
假设某个向量经 $N^{r+2}$ 后为零。

那么它经一次 $N$ 后属于 $\ker N^{r+1}$。由于：

\[
\ker N^{r+1}=\ker N^r,
\]

它经一次 $N$ 后实际上已属于 $\ker N^r$，因此原向量属于 $\ker N^{r+1}$。

所以再高一层不会产生新方向。之后可重复同样推理。
\end{jproof}

\begin{jprop}{有限维空间中的稳定性}{r:10:4}
若：

\[
\dim V=n,
\]

则零空间链最多在第 $n$ 步之前稳定。

因为每次严格增长都会使维数至少增加 $1$，而零空间维数不可能超过 $n$。
\end{jprop}

\begin{jprop}{幂零算子的幂}{r:10:5}
若 $N$ 是 $n$ 维空间上的幂零算子，则：

\[
N^n=0.
\]
\end{jprop}

\begin{jproof}{证明思路}
幂零意味着某一层零空间最终等于整个 $V$。

零空间链至多经过 $n$ 次严格增长便稳定，因此在第 $n$ 步时已经有：

\[
\ker N^n=V.
\]

这等价于：

\[
N^n=0.
\]
\end{jproof}

%% 第 10 章：Fitting 分解
%% 本文件由 tools/md2tex.py 自动生成，请勿直接编辑。

\chapter{Fitting 分解}


\begin{jthm}{Fitting 分解定理}{r:11:1}
设：

\[
A:V\to V
\]

是有限维空间上的算子，且：

\[
n\ge\dim V.
\]

则：

\[
V
=
\ker A^n
\oplus
\operatorname{range}A^n.
\]
\end{jthm}

\section{Fitting 分解的意义}

空间被分为两部分：

\begin{jitem}
\item %
$\ker A^n$：反复作用 $A$ 后最终消失的部分；
\item %
$\operatorname{range}A^n$：经过很多次作用后仍然保留下来的部分。
\end{jitem}

当：

\[
A=T-\lambda I,
\]

第一部分就是 $\lambda$ 的广义特征空间。

\begin{jproof}{证明思路}
当 $n$ 足够大时，零空间链与值域链都已经稳定。

由秩—零化度定理：

\[
\dim\ker A^n+\dim\operatorname{range}A^n
=
\dim V.
\]

因此，只需证明二者交集为零。

若某个向量既在 $\ker A^n$ 中，又在 $\operatorname{range}A^n$ 中，则它可以写成 $A^n(u)$，同时又会被 $A^n$ 消掉。

这意味着 $u$ 落入更高次的零空间。由于零空间链已经稳定，$u$ 其实已经在 $\ker A^n$ 中，因此：

\[
A^n(u)=0.
\]

所以交集只能是零向量。
\end{jproof}

%% 第 11 章：广义特征空间及它的主分解
%% 本文件由 tools/md2tex.py 自动生成，请勿直接编辑。

\chapter{广义特征空间及它的主分解}


\begin{jdefn}{广义特征向量}{r:12:1}
若非零向量 $v$ 满足：

\[
(T-\lambda I)^k v=0
\]

对某个正整数 $k$ 成立，则称 $v$ 是属于 $\lambda$ 的广义特征向量。

普通特征向量是 $k=1$ 时的特殊情形。

类比之前的不同的特征值对应的特征向量线性无关，对于广义特征向量，不同的特征值对应的广义特征向量互相线性无关，证明从略。
\end{jdefn}

\begin{jdefn}{广义特征空间}{r:12:2}
定义：

\[
G_\lambda
=
\ker(T-\lambda I)^{\dim V}.
\]

它称为 $\lambda$ 对应的广义特征空间。

实际上，零空间链可能在更早的某一步稳定；使用 $\dim V$ 只是统一而安全的写法。因为零空间链会递增，并且最远在 $\dim V$ 处停止。
\end{jdefn}

\section{广义特征空间与普通特征空间}

\[
E_\lambda
=
\ker(T-\lambda I)
\subseteq
G_\lambda.
\]

普通特征空间只记录被 $T-\lambda I$ 一次消掉的方向。

广义特征空间则记录经过有限次作用后最终被消掉的所有方向。

\begin{jthm}{不同特征值的广义特征空间直和}{r:12:3}
设：

\[
\lambda_1,\ldots,\lambda_s
\]

是 $T$ 的全部不同特征值。若 $V$ 是有限维复向量空间，则：

\[
V
=
G_{\lambda_1}\oplus\cdots\oplus G_{\lambda_s}.
\]

这称为广义特征空间的主分解。
\end{jthm}

\begin{jproof}{主分解的证明思路}
对第一个特征值 $\lambda_1$，令：

\[
N_1=T-\lambda_1I.
\]

对足够大的 $n$，Fitting 分解给出：

\[
V
=
\ker N_1^n
\oplus
\operatorname{range}N_1^n.
\]

其中：

\[
\ker N_1^n=G_{\lambda_1}.
\]

而：

\[
\operatorname{range}N_1^n
\]

对 $T$ 不变，因此可以将 $T$ 限制到这个剩余空间上继续研究。

在该剩余空间中，$\lambda_1$ 已经不再出现为特征值。对下一个特征值重复这一过程，最终得到全部广义特征空间的直和。

Cayley-Hamilton 定理保证特征多项式中的全部因子作用后会消掉整个空间，因此这个过程最终会覆盖整个 $V$。
\end{jproof}

\section{广义特征空间的维数}

在复数域上：

\[
\dim G_\lambda
=
\lambda\text{ 的代数重数}.
\]

而：

\[
\dim E_\lambda
=
\lambda\text{ 的几何重数}.
\]

代数重数表示广义特征空间的总维数；几何重数表示其中普通特征向量所张成空间的维数。

%% 第 12 章：Jordan 链
%% 本文件由 tools/md2tex.py 自动生成，请勿直接编辑。

\chapter{Jordan 链}


\section{从广义特征空间得到幂零算子}

固定一个特征值 $\lambda$，并限制到其广义特征空间：

\[
G_\lambda.
\]

令：

\[
N=(T-\lambda I)|_{G_\lambda}.
\]

则 $N$ 是 $G_\lambda$ 上的幂零算子。

因此，Jordan 标准型的问题可以转化为：

\begin{jinsight}
如何描述一个幂零算子的内部结构？
\end{jinsight}

\begin{jdefn}{Jordan 链}{r:13:1}
长度为 $r$ 的 Jordan 链是一列向量：

\[
(v_1,\ldots,v_r),
\]

满足：

\[
Nv_1=0,
\]

以及：

\[
Nv_j=v_{j-1},
\qquad
j=2,\ldots,r.
\]

等价地：

\[
Tv_1=\lambda v_1,
\]

\[
Tv_j=\lambda v_j+v_{j-1},
\qquad
j=2,\ldots,r.
\]

其中 $v_1$ 是普通特征向量；后续向量是广义特征向量，但会逐层经 $N$ 回落到前一个向量。
\end{jdefn}

\begin{jprop}{单条 Jordan 链中的向量线性无关}{r:13:2}
若存在向量 $v$，满足：

\[
N^{r}v=0,
\qquad
N^{r-1}v\ne0,
\]

则：

\[
v,Nv,\ldots,N^{r-1}v
\]

线性无关。
\end{jprop}

\begin{jproof}{证明思路}
假设这列向量存在非平凡线性关系。

取其中最低次、且系数不为零的项。对整条关系施加恰当次数的 $N$，所有更高次项都会因 $N^rv=0$ 而消失，只留下该系数乘以：

\[
N^{r-1}v.
\]

但：

\[
N^{r-1}v\ne0.
\]

这与该系数非零矛盾。

因此，一条由链头不断施加 $N$ 生成的完整链，内部一定线性无关。
\end{jproof}

\section{Jordan 链基的构造原则}

一般的幂零算子不一定只需要一条链，也就是说，这个幂零算子的特征值对应的广义特征空间中，可能存在多条完整的链。

不能把：

\[
\ker N,\ker N^2,\ldots
\]

中的所有向量直接全部取来做基，因为这些空间彼此嵌套，其中包含重复方向。

正确做法是：从零空间链的高层开始，挑选真正新增的独立方向作为链头；再沿 $N$ 向下生成链。

设：

\[
K_k=\ker N^k,
\qquad
K_0=\{0\}.
\]

则：

\[
\dim K_k-\dim K_{k-1}
\]

等于长度至少为 $k$ 的 Jordan 链数量。

而：

\[
\bigl(\dim K_k-\dim K_{k-1}\bigr)
-
\bigl(\dim K_{k+1}-\dim K_k\bigr)
\]

等于长度恰好为 $k$ 的 Jordan 链数量。

\begin{jthm}{Jordan 链基存在定理}{r:13:3}
每个有限维幂零算子都存在一组由 Jordan 链组成的基。

在 Jordan 标准型的语境中，先固定一个特征值 $\lambda$，并令：

\[
W=G_\lambda,
\]

\[
N=(T-\lambda I)|_W:W\to W.
\]

此时 $N$ 是 $W$ 上的幂零算子。以下构造的目标是为 $W$ 找到一组 Jordan 链基，而不是直接为原始空间 $V$ 找基。
\end{jthm}

\begin{jidea}{构造的目标}
我们希望选出若干条链：

\[
c,\ Nc,\ N^2c,\ldots
\]

使每条链内部线性无关，并且所有链合在一起后，既不重复，也恰好张成整个 $W$。

这里将 $c$ 称为一条链的生成端：不断对它施加 $N$，便得到整条链。
\end{jidea}

\section{第一步：确定最长链的长度}

取最小的正整数 $r$，使：

\[
N^r=0.
\]

这意味着：对 $W$ 中任何向量连续作用 $N$ 共 $r$ 次，结果都会变为零；但至少存在某个向量，需要恰好 $r$ 次才变成零。

记：

\[
K_k=\ker N^k.
\]

此时：

\[
K_r=W.
\]

因为 $N^r=0$ 表示 $W$ 中每个向量都属于 $\ker N^r$。

\subsection{第二步：先构造最长的链}

$K_{r-1}$ 包含所有经过至多 $r-1$ 次作用就会变为零的向量。

先取 $K_{r-1}$ 的一组基，再将它扩充为 $W$ 的一组基。新补入的向量张成某个子空间，记为 $C_r$：

\[
W=K_{r-1}\oplus C_r.
\]

任取非零 $c\in C_r$。因为：

\[
c\in W=K_r,
\qquad
c\notin K_{r-1},
\]

所以：

\[
N^rc=0,
\qquad
N^{r-1}c\ne0.
\]

因此，$c$ 生成一条长度恰为 $r$ 的链：

\[
c,\ Nc,\ N^2c,\ldots,N^{r-1}c.
\]

对 $C_r$ 的每个基向量都这样做，便得到全部最长的 Jordan 链。

\subsection{第三步：构造较短的链}

接下来处理长度为 $r-1$ 的链。

不能直接从 $K_{r-1}$ 中随便挑向量，因为其中已经包含刚才长链向下作用得到的向量。例如，若：

\[
c,\ Nc,\ldots,N^{r-1}c
\]

是一条长链，那么 $Nc$ 已经位于 $K_{r-1}$ 中；不能再把 $Nc$ 当作另一条新链的生成端。

因此，在第 $k$ 层寻找长度为 $k$ 的链时，要先排除两类已经解释过的方向：

\begin{jitem}
\item %
$K_{k-1}$ 中的向量：它们在不足 $k$ 次作用后已经归零，长度不够；
\item %
已有较长链向下作用后得到的向量：它们已经属于旧链的一部分，不应重复使用。
\end{jitem}

然后，在 $K_k$ 中选取一批无法由这些旧方向表示的新向量。

具体做法仍是熟悉的“扩充基”：

\begin{jenum}
\item %
先取全部旧方向的一组基；
\item %
将它扩充为 $K_k$ 的一组基；
\item %
新补入的向量，就是长度恰为 $k$ 的新链生成端。
\end{jenum}

对每一个新生成端 $c$，依次写出：

\[
c,\ Nc,\ldots,N^{k-1}c.
\]

由于：

\[
c\notin K_{k-1},
\]

它不会提前归零，因此这确实是一条长度为 $k$ 的完整链。

按：

\[
r,r-1,\ldots,1
\]

的顺序重复这个过程，直到所有层都处理完。

\subsection{为什么不同链之间不会线性相关}

单条链内部的线性无关性已经证明过。

不同链之间不会产生冗余，原因在于：每次选择新的生成端时，都先排除了已经属于旧链或能够由旧链表示的方向；随后再通过扩充基，选取真正独立的新方向。

因此，新链不会重复旧链，也不能由已经构造出的链线性表示。

最终，所有链共同张成：

\[
K_1,\ K_2,\ \ldots,\ K_r=W.
\]

也就是说，它们共同张成整个广义特征空间 $W$。

同时，构造过程中每增加一条链，增加的都是此前未出现的独立方向；所有链向量的总数最终恰好等于：

\[
\dim W.
\]

因此，这些向量张成 $W$，且数量恰好等于 $W$ 的维数。它们合起来便构成 $W$ 的一组基。

\subsection{回到原始算子}

对每个特征值 $\lambda$，都在其广义特征空间：

\[
G_\lambda
\]

上重复上述构造。

最后，由主分解：

\[
V=\bigoplus_\lambda G_\lambda,
\]

将所有广义特征空间中的 Jordan 链基合并，便得到原始空间 $V$ 的一组 Jordan 基。

\begin{jinsight}
Jordan 链基的构造不是任意挑选若干条链，而是在每个广义特征空间中，从最长链开始，逐层排除旧方向，再补入新的独立生成端。
\end{jinsight}

%% 第 13 章：Jordan 块与 Jordan 标准型
%% 本文件由 tools/md2tex.py 自动生成，请勿直接编辑。

\chapter{Jordan 块与 Jordan 标准型}


\begin{jdefn}{Jordan 块}{r:14:1}
对长度为 $r$ 的 Jordan 链：

\[
(v_1,\ldots,v_r),
\]

其中：

\[
Tv_1=\lambda v_1,
\]

\[
Tv_j=\lambda v_j+v_{j-1},
\]

则 $T$ 在这组基下的矩阵为：

\[
J_r(\lambda)
=
\begin{pmatrix}
\lambda&1&0&\cdots&0\\
0&\lambda&1&\cdots&0\\
0&0&\lambda&\ddots&0\\
\vdots&\vdots&\ddots&\ddots&1\\
0&0&0&0&\lambda
\end{pmatrix}.
\]

它称为一个属于特征值 $\lambda$ 的 Jordan 块。
\end{jdefn}

\section{Jordan 块的意义}

对角线上的 $\lambda$ 表示该链上的基本缩放。

上方的 $1$ 表示：这个方向除了被缩放外，还会向链中更低一级的方向产生不可消除的耦合。

若块大小大于 $1$，该部分不能仅由普通特征向量组成，因此不能完全对角化。

\begin{jthm}{Jordan 标准型定理}{r:14:2}
若 $V$ 是有限维复向量空间，则任意算子：

\[
T:V\to V
\]

都存在一组基，使其矩阵是若干 Jordan 块构成的块对角矩阵：

\[
[T]
=
J_{r_1}(\lambda_1)
\oplus
\cdots
\oplus
J_{r_m}(\lambda_m).
\]

Jordan 块的排列顺序可以改变，但每个特征值对应的块大小集合是由 $T$ 唯一确定的。
\end{jthm}

\begin{jprop}{Jordan 标准型与对角化}{r:14:3}
\[
T\text{ 可对角化}
\iff
\text{所有 Jordan 块都为 }1\times1.
\]

等价地：

\[
G_\lambda=E_\lambda
\]

对每个特征值 $\lambda$ 都成立。

还等价于最小多项式没有重复根：

\[
m_T(z)
=
\prod_\lambda(z-\lambda).
\]
\end{jprop}

\section{幂零算子的矩阵结构}

若 $N$ 是幂零算子，则它唯一的特征值是：

\[
0.
\]

证明：设 $\lambda$ 是 $N$ 的任意一个特征值，$v\neq 0$ 为对应的特征向量，则

\[
Nv=\lambda v.
\]

两边连续作用算子 $N$，可得

\[
N^2 v = N(Nv)=N(\lambda v)=\lambda Nv=\lambda^2 v.
\]

归纳可得

\[
N^k v=\lambda^k v.
\]

由幂零条件 $N^k=0$，于是

\[
\lambda^k v = N^k v = 0.
\]

因 $v\neq 0$，故 $\lambda^k=0$，从而 $\lambda=0$。

因此幂零算子仅能以 $0$ 作为特征值。

因此，它在 Jordan 基下的矩阵由若干：

\[
J_r(0)
\]

构成。

这些 Jordan 块的对角线全为零。

显然，对于使其成为上三角阵的任意一组基，对角线为零。

幂零算子的矩阵并非在任意基下都显然对角线为零；这里的结论是针对 Jordan 基，或任意使其上三角化的基。

%% 第 14 章：实向量空间的复化
%% 本文件由 tools/md2tex.py 自动生成，请勿直接编辑。

\chapter{实向量空间的复化}


\section{为什么需要复化}

复数域上的有限维算子总能上三角化，并拥有 Jordan 标准型。

实数域上则不一定有实特征值。

例如，平面旋转算子没有非零实特征向量。因此，它无法在实数基下被对角化，也无法在实数基下上三角化。

复化的作用是：把实空间扩展到复数域，使原本隐藏的复特征值与复 Jordan 链显现出来。

\begin{jdefn}{复化空间}{r:15:1}
设 $V$ 是实向量空间。

定义其复化：

\[
V_{\mathbb C}
=
\{x+iy:x,y\in V\}.
\]

它是一个复向量空间。

若：

\[
T:V\to V
\]

是实线性算子，则定义其复化算子：

\[
T_{\mathbb C}:V_{\mathbb C}\to V_{\mathbb C},
\]

\[
T_{\mathbb C}(x+iy)
=
T(x)+iT(y).
\]
\end{jdefn}

\section{复化的矩阵意义}

若 $T$ 在实基下的矩阵是实矩阵 $A$，则 $T_{\mathbb C}$ 在对应复基下的矩阵仍然是同一个 $A$。

变化不在矩阵本身，而在于允许输入、输出坐标取复数。

\begin{jprop}{共轭特征值}{r:15:2}
由于 $T$ 的矩阵系数都是实数，若：

\[
\lambda=a+ib
\]

是 $T_{\mathbb C}$ 的特征值，则：

\[
\overline\lambda=a-ib
\]

也是特征值。

相应的 Jordan 结构也成对出现。
\end{jprop}

%% 第 15 章：复特征值对应的实二维不变子空间
%% 本文件由 tools/md2tex.py 自动生成，请勿直接编辑。

\chapter{复特征值对应的实二维不变子空间}


\section{从复特征向量得到实不变平面}

设：

\[
z=x+iy
\]

是 $T_{\mathbb C}$ 的特征向量，且：

\[
T_{\mathbb C}(z)
=
(a+ib)z,
\qquad
b\ne0.
\]

比较实部与虚部，得到：

\[
T(x)=ax-by,
\]

\[
T(y)=bx+ay.
\]

因此：

\[
\operatorname{span}_{\mathbb R}(x,y)
\]

是 $T$ 的二维实不变子空间。

\section{对应的实矩阵块}

在实基：

\[
(x,y)
\]

下，$T$ 在该二维不变子空间上的矩阵为：

\[
\begin{pmatrix}
a&b\\
-b&a
\end{pmatrix}.
\]

这个二维块同时描述：

\begin{jitem}
\item %
$a$ 所代表的缩放；
\item %
$b$ 所代表的旋转耦合。
\end{jitem}

若 $b=0$，它退化为实特征值对应的一维结构。

\section{复 Jordan 块如何回到实空间}

对于一对共轭复 Jordan 块：

\[
J_r(a+ib),
\qquad
J_r(a-ib),
\]

回到实空间后，它们合并为一个实块矩阵：

\begin{jitem}
\item %
每个复标量 $a+ib$ 替换为：

\[
\begin{pmatrix}
a&b\\
-b&a
\end{pmatrix};
\]
\item %
原 Jordan 块超对角线上的 $1$ 替换为：

\[
I_2.
\]
\end{jitem}

这得到实 Jordan 结构。

\section{复化的意义}

复化让实算子的真实谱结构在更大的数域中显现，再将共轭成对的复结构无损翻译为实空间中的二维不变块。

%% 第 2 部分：内积空间上的算子及其分解
%% 本文件由 tools/md2tex.py 自动生成，请勿直接编辑。

\part{内积空间上的算子及其分解}

\begin{partintro}{这一部分研究什么}
一般向量空间上的算子，可以通过上三角化、对角化与 Jordan 标准型来描述。

内积空间额外拥有长度、正交与投影结构，其中最重要的便是正交。

它允许我们定义伴随算子；当一个算子与其伴随满足特殊关系时，便会获得远强于 Jordan 标准型的谱分解。

\begin{jinsight}
在复内积空间中，正规算子可以在规范正交基下对角化。在实内积空间中，只有自伴算子保证可以在规范正交基下对角化。实正规算子不一定存在实特征向量，例如平面旋转算子。
\end{jinsight}
\end{partintro}

%% 第 16 章：伴随算子
%% 本文件由 tools/md2tex.py 自动生成，请勿直接编辑。

\chapter{伴随算子}


\section{为什么需要定义伴随}

设：

\[
T:V\to W.
\]

若固定 $w\in W$，则：

\[
v\mapsto\langle T(v),w\rangle_W
\]

是 $V$ 上的线性泛函。

它表示：先将 $v$ 通过 $T$ 送到 $W$，再用 $w$ 所代表的内积测量方式对其测量。

\section{伴随的定义}

对固定的 $w\in W$，考虑：

\[
v\mapsto\langle T(v),w\rangle_W.
\]

由于 $T$ 是线性映射，且内积对第一个位置线性，这确实是一个定义在 $V$ 上的线性泛函。

由 Riesz 表示定理，$V$ 上的每个线性泛函都能被唯一写成“与 $V$ 中某个向量做内积”的形式。因此，存在唯一的向量 $T^*(w)\in V$，使：

\[
\langle T(v),w\rangle_W
=
\langle v,T^*(w)\rangle_V.
\]

这里，$T^*(w)$ 就是将 $W$ 中由 $w$ 给出的测量方式，经由 $T$ 拉回到 $V$ 后所对应的唯一向量。

伴随与此前的对偶映射本质上都是对于泛函这种测量方式在输入输出空间中的转换。

但区别在于：

\begin{jitem}
\item %
对偶映射将 $W$ 中的泛函拉回为 $V$ 中的泛函；也就是说，$T'(\varphi)=\varphi\circ T.$
\item %
伴随利用 Riesz 表示定理，将这个被拉回去的泛函，直接用向量而不是复合的映射的方式，在 $V$ 中表示出来，并且表示为一种内积的形式。
\end{jitem}

\begin{jdefn}{伴随算子的定义}{r:17:1}
伴随算子是唯一满足下式的线性映射：

\[
T^*:W\to V,
\]

\[
\langle T(v),w\rangle_W
=
\langle v,T^*(w)\rangle_V.
\]
\end{jdefn}

\begin{jprop}{伴随算子的存在与唯一性}{r:17:2}
对于每个 $w\in W$，映射：

\[
v\mapsto\langle T(v),w\rangle_W
\]

是 $V$ 上的线性泛函。

Riesz 表示定理保证存在唯一向量 $T^*(w)$ 表示这个泛函，因此伴随存在且唯一。

$T^*$ 具有加性和齐性，因此 $T^*$ 是线性映射。
\end{jprop}

\begin{jprop}{伴随的代数性质}{r:17:3}
设：

\[
S,T:V\to W,
\]

且 $\lambda\in\mathbb F$。

则：

\[
(S+T)^*=S^*+T^*,
\]

\[
(\lambda T)^*=\overline{\lambda}T^*,
\]

\[
(T^*)^*=T.
\]

若：

\[
T:U\to V,
\qquad
S:V\to W,
\]

则：

\[
(S\circ T)^*=T^*\circ S^*.
\]
\end{jprop}

\begin{jproof}{这些性质的证明思路}
证明使用同一种思路：将等式两边放进内积，并不断使用伴随定义把算子从内积的第一个位置移到第二个位置。

加法对应内积的加法性。

标量取共轭，是因为本笔记采用第一位置线性的内积约定：标量从第二个位置移出时会变为共轭。

复合反序只需要通过反复应用伴随的定义在内积的表达式上即可。
\end{jproof}

%% 第 17 章：伴随的零空间和值域
%% 本文件由 tools/md2tex.py 自动生成，请勿直接编辑。

\chapter{伴随的零空间和值域}


\begin{jprop}{伴随的零空间}{r:18:1}
\[
\ker T^*
=
(\operatorname{range}T)^\perp.
\]
\end{jprop}

\begin{jproof}{证明思路}
$w\in\ker T^*$，意味着：

\[
T^*(w)=0.
\]

由伴随定义，这等价于：

\[
\langle T(v),w\rangle=0
\]

对所有 $v\in V$ 都成立。

而所有 $T(v)$ 恰好组成 $\operatorname{range}T$。因此，$w$ 与 $T$ 的全部输出正交，即：

\[
w\in(\operatorname{range}T)^\perp.
\]
\end{jproof}

\begin{jprop}{伴随的值域}{r:18:2}
在有限维情形：

\[
\operatorname{range}T^*
=
(\ker T)^\perp.
\]
\end{jprop}

\begin{jproof}{证明思路}
若 $u=T^*(w)$，且 $v\in\ker T$，则：

\[
\langle v,u\rangle
=
\langle v,T^*(w)\rangle
=
\langle T(v),w\rangle
=
0.
\]

因此：

\[
\operatorname{range}T^*
\subseteq
(\ker T)^\perp.
\]

另一方面，伴随与原映射具有相同的秩：

\[
\operatorname{rank}T^*=\operatorname{rank}T.
\]

而：

\[
\dim(\ker T)^\perp
=
\dim V-\dim\ker T
=
\operatorname{rank}T.
\]

一个子空间包含于另一个子空间，且两者维数相同，因此它们相等。
\end{jproof}

\section{四个子空间关系}

伴随将此前的四个子空间关系写成：

\[
\ker T^*
=
(\operatorname{range}T)^\perp,
\]

\[
\operatorname{range}T^*
=
(\ker T)^\perp.
\]

在标准欧氏坐标下，$T^*$ 对应共轭转置；实矩阵情形下就是转置。

这正是经典“四个基本子空间”中行空间、列空间、零空间与左零空间关系的抽象版本，我们将在表示的笔记中更加详细的去说明。

%% 第 18 章：伴随的矩阵
%% 本文件由 tools/md2tex.py 自动生成，请勿直接编辑。

\chapter{伴随的矩阵}


\begin{jprop}{规范正交基下的共轭转置}{r:19:1}
若 $V,W$ 分别选取规范正交基，且：

\[
[T]=A,
\]

则：

\[
[T^*]=A^*.
\]

这里：

\[
A^*={A}^{\mathsf H}
\]

称为 $A$ 的共轭转置。

在实数域上：

\[
A^*=A^{\mathsf T}.
\]
\end{jprop}

\section{为什么必须使用规范正交基}

只有在规范正交基下，坐标中的内积才直接等于普通欧氏点积。

若基不是规范正交基，坐标向量之间的普通点积并不等于原空间中的内积。此时必须先通过 Gram 矩阵记录该基的内积关系，再进行相应转换。

因此，伴随矩阵等于共轭转置”是规范正交坐标下的结论

%% 第 19 章：自伴算子
%% 本文件由 tools/md2tex.py 自动生成，请勿直接编辑。

\chapter{自伴算子}


\begin{jdefn}{自伴算子的定义}{r:20:1}
若：

\[
T=T^*,
\]

则称 $T$ 是自伴算子。

在规范正交基下，自伴算子的矩阵满足：

\[
A=A^*.
\]

实数情形中，这就是对称矩阵；复数情形中，这就是 Hermitian 矩阵（共轭对称阵）。
\end{jdefn}

\begin{jprop}{自伴算子的基本性质}{r:20:2}
若 $S,T$ 都是自伴算子，则：

\[
S+T
\]

也是自伴算子。

若 $a\in\mathbb R$，则：

\[
aT
\]

也是自伴算子。

两个自伴算子的乘积不一定自伴。若额外满足：

\[
ST=TS,
\]

则：

\[
(ST)^*=T^*S^*=TS=ST,
\]

所以 $ST$ 自伴。
\end{jprop}

\begin{jprop}{自伴算子的特征值是实数}{r:20:3}
若：

\[
Tv=\lambda v,
\qquad
v\ne0,
\]

且 $T$ 自伴，则：

\[
\lambda\in\mathbb R.
\]
\end{jprop}

\begin{jproof}{证明思路}
计算：

\[
\langle Tv,v\rangle.
\]

一方面，由特征值关系，它等于 $\lambda\langle v,v\rangle$。

另一方面，由自伴性：

\[
\langle Tv,v\rangle
=
\langle v,Tv\rangle,
\]

所以它等于自己的共轭。

由于：

\[
\langle v,v\rangle>0,
\]

只能有：

\[
\lambda=\overline{\lambda}.
\]

因此 $\lambda$ 是实数。
\end{jproof}

\begin{jprop}{复内积空间中自伴的刻画}{r:20:4}
在复内积空间中：

\[
T\text{ 自伴}
\iff
\langle Tv,v\rangle\in\mathbb R
\text{ 对所有 }v\in V.
\]
\end{jprop}

\begin{jproof}{证明思路}
自伴时，内积可以直接看出等于自己的共轭，因此为实数。

反过来，若所有：

\[
\langle Tv,v\rangle
\]

都为实数，则分别对 $u+v$ 与 $u+iv$ 使用该条件。

展开后可分别得到：

\[
\langle Tu,v\rangle
=
\langle u,Tv\rangle.
\]

这正是：

\[
T=T^*.
\]
\end{jproof}

\begin{jprop}{复内积空间中的零二次型结论}{r:20:5}
若 $V$ 是复内积空间，且：

\[
\langle Tv,v\rangle=0
\]

对所有 $v\in V$ 成立，则：

\[
T=0.
\]
\end{jprop}

\begin{jproof}{证明思路}
分别代入 $u+v$ 与 $u+iv$。

前者表明：

\[
\langle Tu,v\rangle
\]

的实部为零；后者表明其虚部也为零。

因此：

\[
\langle Tu,v\rangle=0
\]

对所有 $u,v$ 成立。

特别取：

\[
v=Tu,
\]

得到：

\[
\|Tu\|^2=0.
\]

所以 $Tu=0$ 对所有 $u$ 成立，即：

\[
T=0.
\]

实内积空间中此结论不成立。平面上的 $90^\circ$ 旋转算子非零，但对每个 $v$ 都满足：

\[
\langle Tv,v\rangle=0.
\]
\end{jproof}

%% 第 20 章：正规算子
%% 本文件由 tools/md2tex.py 自动生成，请勿直接编辑。

\chapter{正规算子}


\begin{jdefn}{正规算子的定义}{r:21:1}
若：

\[
T^*T=TT^*,
\]

则称 $T$ 是正规算子。

自伴算子一定正规，因为：

\[
T^*T=TT=TT^*.
\]

但正规算子不一定自伴。平面旋转算子就是正规但非自伴的例子。
\end{jdefn}

\begin{jprop}{正规算子的范数刻画}{r:21:2}
\[
T\text{ 正规}
\iff
\|Tv\|=\|T^*v\|
\text{ 对所有 }v\in V.
\]
\end{jprop}

\begin{jproof}{证明思路}
将两边范数平方展开：

\[
\|Tv\|^2
=
\langle T^*Tv,v\rangle,
\]

\[
\|T^*v\|^2
=
\langle TT^*v,v\rangle.
\]

因此，若 $T^*T=TT^*$，两者范数相同。

反过来，若所有向量的两个范数都相同，在复内积空间中，这足以推出：

\[
T^*T=TT^*.
\]
\end{jproof}

\begin{jprop}{正规算子的特征向量也是伴随的特征向量}{r:21:3}
若 $T$ 正规，且：

\[
Tv=\lambda v,
\qquad
v\ne0,
\]

则：

\[
T^*v=\overline{\lambda}v.
\]
\end{jprop}

\begin{jproof}{证明思路}
算子：

\[
T-\lambda I
\]

仍然正规。

而：

\[
(T-\lambda I)v=0.
\]

正规算子的范数刻画说明：若一个向量被正规算子映为零，它也会被该算子的伴随映为零。

因此：

\[
(T-\lambda I)^*v=0.
\]

又因为：

\[
(T-\lambda I)^*
=
T^*-\overline{\lambda}I,
\]

所以：

\[
T^*v=\overline{\lambda}v.
\]
\end{jproof}

\begin{jprop}{不同特征值对应的特征向量正交}{r:21:4}
若 $T$ 正规，且：

\[
Tv=\lambda v,
\qquad
Tw=\mu w,
\qquad
\lambda\ne\mu,
\]

则：

\[
\langle v,w\rangle=0.
\]
\end{jprop}

\begin{jproof}{证明思路}
由前一结论：

\[
T^*w=\overline{\mu}w.
\]

于是：

\[
\langle Tv,w\rangle
=
\lambda\langle v,w\rangle.
\]

另一方面：

\[
\langle Tv,w\rangle
=
\langle v,T^*w\rangle
=
\mu\langle v,w\rangle.
\]

因此：

\[
(\lambda-\mu)\langle v,w\rangle=0.
\]

因为：

\[
\lambda\ne\mu,
\]

所以：

\[
\langle v,w\rangle=0.
\]
\end{jproof}

%% 第 21 章：复谱定理
%% 本文件由 tools/md2tex.py 自动生成，请勿直接编辑。

\chapter{复谱定理}


\begin{jthm}{复谱定理}{r:22:1}
设 $V$ 是有限维复内积空间，$T:V\to V$。

下列条件等价：

\begin{jenum}
\item %
$T$ 是正规算子；
\item %
$V$ 存在一组由 $T$ 的特征向量构成的规范正交基；
\item %
$T$ 在某组规范正交基下具有对角矩阵。
\end{jenum}
\end{jthm}

\section{复谱定理的意义}

一般复算子只能保证上三角化，Jordan 块仍可能表示方向之间的耦合。

正规性保证这些耦合可以完全消失。空间可以由两两正交的特征方向组成，算子只在每个方向上独立地乘以一个复数特征值。

因此，正规算子在内积空间下是最理想的可对角化算子。

\begin{jproof}{复谱定理的证明思路}
Schur 定理保证：任意复内积空间上的算子，都能在某组规范正交基下表示为上三角矩阵。

设该上三角矩阵为 $A$。

若 $T$ 正规，则：

\[
\|Ae_1\|=\|A^*e_1\|.
\]

对于上三角矩阵，$Ae_1$ 只包含第一个对角元的贡献；而 $A^*e_1$ 还会包含第一行其余元素的贡献。

两个范数相等，只能说明第一行除对角元外的全部元素都为零。

于是 $A$ 被分成一个 $1\times1$ 块和一个更小的正规上三角块。对余下的块重复同样论证，最终所有非对角元素都为零。

因此，$A$ 必为对角矩阵。
\end{jproof}

\section{谱定理中的对角元}

若规范正交特征基为：

\[
(e_1,\ldots,e_n),
\]

且：

\[
Te_j=\lambda_je_j,
\]

则：

\[
[T]
=
\operatorname{diag}(\lambda_1,\ldots,\lambda_n).
\]

对角元就是对应基向量的特征值，按代数重数重复出现。

谱定理只能保证这种对角矩阵存在，但不给出每个特征值的具体数值，也就是说，谱定理保证对于这样的算子来说，存在某组正交基，且这组正交基的是整个空间的直和项，且每一项都是一维不变子空间，也就是一个特征向量。具体的特征值需要我们自己计算。

更一般的说法是：谱定理保证了这个算子在一个由它的特征向量构成的正交基下存在一个对角阵。

\section{规范正交特征基不唯一}

若某个特征值的特征空间维数大于 $1$，可以在该特征空间内部任取规范正交基。

不同特征值对应的特征空间彼此正交，因此将各特征空间中的规范正交基合并，仍得到整个空间的规范正交特征基。

即使全部特征值互异，也可以重排基向量；在复数域中，还可以分别给每个特征向量乘以模为 $1$ 的复数。

因此，谱定理通常给出许多不同的规范正交特征基。

%% 第 22 章：实谱定理
%% 本文件由 tools/md2tex.py 自动生成，请勿直接编辑。

\chapter{实谱定理}


\begin{jthm}{实谱定理}{r:23:1}
设 $V$ 是有限维实内积空间，$T:V\to V$。

下列条件等价：

\begin{jenum}
\item %
$T$ 是自伴算子；
\item %
$V$ 存在一组由 $T$ 的特征向量构成的规范正交基；
\item %
$T$ 在某组规范正交基下具有对角矩阵。
\end{jenum}

这里对角矩阵的全部对角元都是实特征值。
\end{jthm}

\section{为什么实谱定理要求自伴而非正规}

自伴算子一定正规。

但实正规算子不一定有实特征值。例如二维旋转算子保范且正规，却没有非零实特征向量，因此不能在实规范正交基下对角化。

自伴性比正规性更强，它排除了旋转型的二维耦合，只保留沿实特征方向的伸缩。

\begin{jprop}{一个二次算子可逆}{r:23:2}
若 $T$ 自伴，$a\in\mathbb R$，且 $b\in\mathbb R$ 非零，则：

\[
(T-aI)^2+b^2I
\]

可逆。
\end{jprop}

\begin{jproof}{证明思路}
$T-aI$ 仍然自伴。

对任意非零 $v$：

\[
\left\langle
\bigl((T-aI)^2+b^2I\bigr)v,v
\right\rangle
=
\|(T-aI)v\|^2+b^2\|v\|^2.
\]

右边严格大于零。

所以该算子的核只有零向量。在有限维空间中，单射等价于可逆，因此它可逆。
\end{jproof}

\section{实自伴算子一定有实特征值}


\begin{jproof}{证明思路}
先将 $T$ 复化，得到复内积空间上的算子 $T_{\mathbb C}$。

复数域上，$T_{\mathbb C}$ 一定有特征值：

\[
\lambda=a+ib.
\]

若：

\[
b\ne0,
\]

则复化后的算子：

\[
(T_{\mathbb C}-aI)^2+b^2I
\]

会将对应特征向量映为零。

但它正是前一个可逆算子的复化，因此不可能有非零核，产生矛盾。

所以：

\[
b=0.
\]

于是 $T$ 存在实特征值。
\end{jproof}

\begin{jprop}{自伴算子与不变子空间}{r:23:3}
设 $T$ 自伴，且 $U$ 是 $T$ 的不变子空间。

则：

\[
U^\perp
\]

也是 $T$ 的不变子空间。
\end{jprop}

\begin{jproof}{证明思路}
任取：

\[
w\in U^\perp,
\qquad
u\in U.
\]

因为 $U$ 对 $T$ 不变，所以：

\[
T(u)\in U.
\]

由自伴性：

\[
\langle T(w),u\rangle
=
\langle w,T(u)\rangle
=
0.
\]

因此：

\[
T(w)\in U^\perp.
\]
\end{jproof}

\begin{jprop}{限制算子仍然自伴}{r:23:4}
在上述条件下：

\[
T|_U:U\to U
\]

与：

\[
T|_{U^\perp}:U^\perp\to U^\perp
\]

都是自伴算子。
\end{jprop}

\begin{jproof}{证明思路}
自伴定义中的内积等式，对 $U$ 或 $U^\perp$ 内的向量仍然成立；而两者又各自不变，所以限制后的算子确实仍映回相应子空间。
\end{jproof}

\begin{jproof}{实谱定理的证明思路}
由前置引理，$T$ 存在一个单位特征向量 $e_1$。

于是：

\[
\operatorname{span}(e_1)
\]

是一维不变子空间。

其正交补：

\[
\operatorname{span}(e_1)^\perp
\]

也对 $T$ 不变，且限制算子仍自伴。

在这个维数减少一的正交补上重复同样过程。最终得到：

\[
(e_1,\ldots,e_n)
\]

这组规范正交特征基。
\end{jproof}

%% 第 23 章：正算子与正平方根
%% 本文件由 tools/md2tex.py 自动生成，请勿直接编辑。

\chapter{正算子与正平方根}


\begin{jdefn}{正算子的定义}{r:24:1}
若 $T$ 自伴，且：

\[
\langle Tv,v\rangle\ge0
\]

对所有 $v\in V$ 成立，则称 $T$ 是正算子。

正算子描述一种不会在任何方向上产生负二次能量的变换。这便于我们所谓的二次型产生联系了。
\end{jdefn}

\begin{jprop}{正算子的谱刻画}{r:24:2}
\[
T\text{ 正}
\iff
T\text{ 自伴，且所有特征值非负}.
\]
\end{jprop}

\begin{jproof}{证明思路}
若：

\[
Te_j=\lambda_je_j,
\]

则：

\[
\langle Te_j,e_j\rangle
=
\lambda_j.
\]

所以正性要求每个特征值都非负。

反过来，在规范正交特征基下：

\[
v=\sum_j a_je_j.
\]

于是：

\[
\langle Tv,v\rangle
=
\sum_j\lambda_j|a_j|^2.
\]

若全部 $\lambda_j\ge0$，该和必然非负。

又因为前面证明过，在 $\mathbb C^n$ 上只有自伴算子才能使二次型都是实数，所以，正算子是自伴的。
\end{jproof}

\begin{jthm}{正平方根}{r:24:3}
若 $T$ 是正算子，则存在唯一的正算子 $R$，使：

\[
R^2=T.
\]

记作：

\[
R=T^{1/2}.
\]
\end{jthm}

\section{正平方根的构造}

在 $T$ 的规范正交特征基下：

\[
Te_j=\lambda_je_j,
\qquad
\lambda_j\ge0.
\]

定义：

\[
Re_j=\sqrt{\lambda_j}e_j.
\]

则：

\[
R^2e_j=\lambda_je_j=Te_j.
\]

因此：

\[
R^2=T.
\]

\section{正平方根的唯一性}

正平方根在每个特征方向上的特征值必须是非负数：

\[
\sqrt{\lambda_j}.
\]

所以它被 $T$ 的特征结构唯一确定。（由前面的定理，线性映射在给定在一组基上的取值时候，这个映射便被唯一确定了)

注意：普通平方根未必唯一。例如恒等算子在不同特征方向上可以取 $1$ 或 $-1$；但其中只有全取正的那个平方根是正平方根。

%% 第 24 章：等距同构、酉算子与正交算子
%% 本文件由 tools/md2tex.py 自动生成，请勿直接编辑。

\chapter{等距同构、酉算子与正交算子}


\begin{jdefn}{等距映射}{r:25:1}
线性映射：

\[
S:V\to W
\]

若满足：

\[
\|S(v)\|=\|v\|
\]

对所有 $v\in V$ 成立，则称 $S$ 是等距映射。

在线性映射情形，保范等价于保内积：

\[
\langle S(u),S(v)\rangle_W
=
\langle u,v\rangle_V.
\]

这可由实数或复数的极化公式推出。
\end{jdefn}

\begin{jdefn}{等距同构}{r:25:2}
若 $S$ 既是等距映射又是双射，则称 $S$ 是等距同构。

当定义域与陪域是同一内积空间时：

\begin{jitem}
\item %
实数域上的等距同构称为正交算子；
\item %
复数域上的等距同构称为酉算子。
\end{jitem}
\end{jdefn}

\begin{jprop}{等距同构与伴随}{r:25:3}
若：

\[
S:V\to W
\]

是等距同构，则：

\[
S^*=S^{-1}.
\]

因此：

\[
S^*S=I_V,
\qquad
SS^*=I_W.
\]

若 $S:V\to V$，则它是正规算子：

\[
S^*S=SS^*=I.
\]

等距同构一般不自伴。

例如平面上的 $90^\circ$ 旋转是正交算子，但：

\[
S^*\ne S.
\]

只有额外满足：

\[
S=S^*
\]

时，等距同构才同时是自伴的。
\end{jprop}

%% 第 3 部分：对于更一般的线性映射的分解
%% 本文件由 tools/md2tex.py 自动生成，请勿直接编辑。

\part{对于更一般的线性映射的分解}

\begin{partintro}{这一部分研究什么}
极分解与奇异值分解更进一步：它们不要求映射从空间映到自身，因此可用于一般线性映射：

\[
T:V\to W.
\]

我们通过前面研究的这些特殊的算子，尝试通过这些算子找出一般的线性映射，是否具有简单的直观的分解。
\end{partintro}

%% 第 25 章：极分解
%% 本文件由 tools/md2tex.py 自动生成，请勿直接编辑。

\chapter{极分解}


\section{为什么需要极分解}

谱定理只处理：

\[
T:V\to V.
\]

但大量线性映射的定义域与陪域不同：

\[
T:V\to W.
\]

极分解将一般映射拆成两件事：

\begin{jitem}
\item %
输入方向被拉伸多少；
\item %
拉伸后如何以保长度的方式搬运到输出空间。
\end{jitem}

\begin{jdefn}{绝对值算子}{r:26:1}
对：

\[
T:V\to W,
\]

算子：

\[
T^*T:V\to V
\]

是正算子，因为：

\[
\langle T^*Tv,v\rangle
=
\langle Tv,Tv\rangle
=
\|Tv\|^2
\ge0.
\]

因此可定义：

\[
|T|
=
(T^*T)^{1/2}.
\]

它称为 $T$ 的绝对值算子。
\end{jdefn}

\begin{jthm}{极分解定理}{r:26:2}
存在唯一的部分等距算子：

\[
U:V\to W,
\]

使：

\[
T=U|T|.
\]

其中：

\[
|T|=(T^*T)^{1/2}.
\]
\end{jthm}

\section{极分解的意义}

$|T|$ 是定义在输入空间上的正算子，它描述各输入方向被放大的程度。

$U$ 则负责方向的搬运。

更准确地说，$U$ 在：

\[
(\ker T)^\perp
\]

到：

\[
\operatorname{range}T
\]

之间是等距同构；而在：

\[
\ker T
\]

上为零。

若 $T$ 可逆，则：

\[
U
\]

是整个空间上的酉算子或正交算子。

%% 第 26 章：奇异值分解
%% 本文件由 tools/md2tex.py 自动生成，请勿直接编辑。

\chapter{奇异值分解}


\section{为什么需要奇异值分解}

特征值分解要求：

\[
T:V\to V.
\]

而奇异值分解适用于任意有限维内积空间之间的线性映射：

\[
T:V\to W.
\]

它同时在输入空间和输出空间中选择规范正交基，使 $T$ 只沿对应方向进行独立缩放。

\section[从 T*T 出发]{从 $T^*T$ 出发}

算子：

\[
T^*T:V\to V
\]

是自伴正算子。

由谱定理，存在 $V$ 的规范正交基：

\[
(v_1,\ldots,v_n),
\]

使：

\[
T^*Tv_i=\sigma_i^2v_i,
\qquad
\sigma_i\ge0.
\]

其中非负数：

\[
\sigma_i
\]

称为 $T$ 的奇异值。

\begin{jprop}{非零奇异值对应的输出方向}{r:27:1}
设：

\[
\sigma_i>0.
\]

定义：

\[
u_i
=
\frac{T(v_i)}{\sigma_i}.
\]

则：

\[
T(v_i)=\sigma_iu_i.
\]

并且全部这样的 $u_i$ 构成 $W$ 中的规范正交列表。
\end{jprop}

\begin{jproof}{证明思路}
对任意正奇异值对应的 $v_i,v_j$：

\[
\langle T(v_i),T(v_j)\rangle
=
\langle v_i,T^*Tv_j\rangle.
\]

而：

\[
T^*Tv_j=\sigma_j^2v_j.
\]

因为 $v_i$ 彼此规范正交，得到：

\[
\langle T(v_i),T(v_j)\rangle
=
\sigma_i^2\delta_{ij}.
\]

除以对应的 $\sigma_i,\sigma_j$ 后，便得到：

\[
\langle u_i,u_j\rangle=\delta_{ij}.
\]
\end{jproof}

\begin{jprop}{零奇异值与零空间}{r:27:2}
若：

\[
\sigma_i=0,
\]

则：

\[
T^*Tv_i=0.
\]

于是：

\[
\|Tv_i\|^2
=
\langle T^*Tv_i,v_i\rangle
=
0.
\]

因此：

\[
Tv_i=0.
\]

所以，零奇异值对应的右奇异向量恰好张成：

\[
\ker T.
\]

这就是 SVD 中“降维”发生的位置：这些输入方向会被完全压缩掉。
\end{jprop}

\begin{jthm}{奇异值分解定理}{r:27:3}
设：

\[
T:V\to W
\]

是有限维内积空间之间的线性映射，且：

\[
\operatorname{rank}T=r.
\]

则存在 $V$ 的规范正交基：

\[
(v_1,\ldots,v_n),
\]

以及 $W$ 的规范正交基：

\[
(u_1,\ldots,u_m),
\]

使：

\[
T(v_i)=\sigma_i u_i,
\qquad
i=1,\ldots,r,
\]

其中：

\[
\sigma_1,\ldots,\sigma_r>0.
\]

并且：

\[
T(v_i)=0,
\qquad
i=r+1,\ldots,n.
\]

在这两组基下，$T$ 的矩阵为矩形对角矩阵：

\[
\Sigma
=
\begin{pmatrix}
\sigma_1 & 0 & \cdots & 0\\
0 & \sigma_2 & \cdots & 0\\
\vdots & \vdots & \ddots & \vdots\\
0 & 0 & \cdots & \sigma_r\\
0 & 0 & \cdots & 0\\
\vdots & \vdots & & \vdots
\end{pmatrix}.
\]

矩阵形式常写为：

\[
A=U\Sigma V^*.
\]
\end{jthm}

\section{SVD 为什么一定有足够多的基向量}

第一步，谱定理为：

\[
T^*T
\]

提供 $V$ 的完整规范正交特征基。

其中正特征值方向经过 $T$ 后，构造出 $W$ 中规范正交的左奇异向量。

零特征值方向正好组成：

\[
\ker T.
\]

最后，将已有左奇异向量补全为 $W$ 的规范正交基。

因此，SVD 不假设输入输出空间恰好存在配对基（这是我们先前学习的一大误解），它由：

\[
T^*T
\]

的正谱结构主动构造输入基，再由 $T$ 构造输出基。

\section[从 TT* 出发]{从 $TT^*$ 出发}

同样可以从：

\[
TT^*:W\to W
\]

开始。

若：

\[
T^*Tv=\sigma^2v,
\qquad
\sigma>0,
\]

则：

\[
TT^*(Tv)
=
\sigma^2(Tv).
\]

因此，$T^*T$ 与 $TT^*$ 的全部正特征值相同，且重数相同。

它们分别描述同一组奇异模式在输入空间与输出空间中的表现：

\begin{jitem}
\item %
$T^*T$ 给出右奇异向量，即输入方向；
\item %
$TT^*$ 给出左奇异向量，即输出方向。
\end{jitem}

二者在零特征值部分可能维数不同，这正反映了定义域、陪域维数与核、余核结构的差异。



% =========================================================================
%  后置部分：定理与定义一览
% =========================================================================
\backmatter
\cleardoublepage
\phantomsection
\addcontentsline{toc}{chapter}{定义、定理与命题一览}
\markboth{定义、定理与命题一览}{定义、定理与命题一览}
\chapter*{定义、定理与命题一览}
\tcblistof[\section*]{jdeflist}{定义}
\tcblistof[\section*]{jthmlist}{定理}
\tcblistof[\section*]{jproplist}{命题}

\end{document}
